% Figure: the mechanism of a cascading transmission overload. A generator
% feeds a load over three parallel corridors. When the first line trips, its
% flow redistributes onto the survivors; if they are already near their
% thermal limit the next overloads and trips, and the redistribution repeats
% until the corridor collapses. The three panels are the same network at
% three moments: intact, one line lost, two lines lost.
\begin{figure}[htbp]
\centering
\begin{tikzpicture}[
  gen/.style={circle, draw=engblue, very thick, minimum size=10mm, font=\scriptsize},
  load/.style={rectangle, draw=engred, very thick, minimum size=8mm, font=\scriptsize},
  intact/.style={engblue, thick},
  tripped/.style={enggray, thick, dash pattern=on 1pt off 2pt},
  hot/.style={engorange, very thick},
  lbl/.style={font=\scriptsize\itshape, color=enggray},
]
% ----- panel 1: intact -----
\begin{scope}[xshift=0cm]
  \node[gen] (g1) at (0,0) {G};
  \node[load] (l1) at (3.4,0) {L};
  \draw[intact] (g1) to[out=35,in=145] node[above,font=\scriptsize]{$1.0$} (l1);
  \draw[intact] (g1) -- node[above,font=\scriptsize]{$1.0$} (l1);
  \draw[intact] (g1) to[out=-35,in=-145] node[below,font=\scriptsize]{$1.0$} (l1);
  \node[lbl] at (1.7,-1.6) {intact: 3 lines};
\end{scope}
% ----- panel 2: one line lost -----
\begin{scope}[xshift=5.6cm]
  \node[gen] (g2) at (0,0) {G};
  \node[load] (l2) at (3.4,0) {L};
  \draw[tripped] (g2) to[out=35,in=145] node[above,font=\scriptsize]{trip} (l2);
  \draw[hot] (g2) -- node[above,font=\scriptsize]{$1.5$} (l2);
  \draw[hot] (g2) to[out=-35,in=-145] node[below,font=\scriptsize]{$1.5$} (l2);
  \node[lbl] at (1.7,-1.6) {overload on survivors};
\end{scope}
% ----- panel 3: two lines lost -----
\begin{scope}[xshift=11.2cm]
  \node[gen] (g3) at (0,0) {G};
  \node[load] (l3) at (3.4,0) {L};
  \draw[tripped] (g3) to[out=35,in=145] (l3);
  \draw[tripped] (g3) -- (l3);
  \draw[tripped] (g3) to[out=-35,in=-145] node[below,font=\scriptsize]{trip} (l3);
  \node[lbl] at (1.7,-1.6) {corridor collapse};
\end{scope}
\end{tikzpicture}
\caption[The cascading-overload mechanism]{The mechanism of a cascading
transmission overload, the same corridor at three moments. Left: a generator
$\mathrm{G}$ feeds a load $\mathrm{L}$ over three parallel lines, each
carrying its share (per-unit flow shown). Centre: the top line trips, and its
flow redistributes onto the two survivors, which now carry $1.5$ times their
former load and run hot (orange), past their thermal rating. Right: an
overloaded survivor trips in turn, the flow redistributes again onto the last
line, which cannot carry the whole load, and the corridor collapses. Each
redistribution is faster than the last, which is why a cascade outruns
operator intervention once it starts and why the $N-1$ contingency standard
requires the survivors to hold the full flow after any single loss.}
\label{fig:vol04ch14:cascade-network}
\end{figure}
